	%----------------------------------------------------------------------------------------
%	TITLE PAGE
%----------------------------------------------------------------------------------------
\KOMAoptions{fontsize=12pt} 
\begin{titlepage} % Suppresses displaying the page number on the title page and the subsequent page counts as page 1
	\newcommand{\HRule}{\rule{\linewidth}{0.5mm}} % Defines a new command for horizontal lines, change thickness here
	
	\centering % Centre everything on the page
	
	\Large МІНІСТЕРСТВО ОСВІТИ І НАУКИ УКРАЇНИ\\  % Main heading such as the name of your university/college
	\medskip
	\large ОДЕСЬКИЙ НАЦІОНАЛЬНИЙ УНІВЕРСИТЕТ ІМЕНІ І.І. МЕЧНИКОВА% Main heading such as the name of your university/college
	
	\large ФАКУЛЬТЕТ МАТЕМАТИКИ, ФІЗИКИ ТА ІНФОРМАЦІЙНИХ ТЕХНОЛОГІЙ % Major heading such as course name
	
	%\textbf{\Large Кафедра Оптимального Керування та Економічної Кібернетики}\\[1.5cm] % Minor heading such as course title
	
	\vspace*{5\baselineskip} % Whitespace above the title
	
	
	{\Large 	\textbf{Євгенія ПЛАТОНОВА, Євген СТРАХОВ}\\} % Editor list
	
	\vspace{1\baselineskip} % Whitespace below the editor list
	\rule{\textwidth}{1.6pt}\vspace*{-\baselineskip}\vspace*{2pt} % Thick horizontal rule
	\rule{\textwidth}{0.4pt} % Thin horizontal rule
	
	\vspace{1\baselineskip} % Whitespace above the title
	
	{\LARGE \textbf{Програмування}} % Title
	
	\vspace{0.5\baselineskip} % Whitespace below the title
	
	\rule{\textwidth}{0.4pt}\vspace*{-\baselineskip}\vspace{3.2pt} % Thin horizontal rule
	\rule{\textwidth}{1.6pt} % Thick horizontal rule
	
	\vspace{1\baselineskip} % Whitespace after the title block
	
	%------------------------------------------------
	%	Subtitle
	%------------------------------------------------
	
	\LARGE \textbf{Основи програмування на мові Python.\\ Структурне програмування\\} % Subtitle or further description
	%{(адаптація learnopengl.com)\\}
	
	\vspace*{0.5\baselineskip} % Whitespace under the subtitle
	
	\Large \textbf{Конспект лекцій}
	\vspace*{1.5\baselineskip}
	%\vfill\v% Position the date 3/4 down the remaining page
	
	%		Рекомендовано Методичною радою ОНУ імені І.І. Мечникова\\
	%		як навчальний посібник для здобувачів ступеня бакалавра\\ 
	%		спеціальності 113 «Прикладна математика», 122 «Комп’ютерні науки»,\\
	%		123 «Комп’ютерна інженерія», 126 «Інформаційні системи та технології»\\
	
	
	\vspace*{0.25\baselineskip}
	
	%	\texttt{Електронне мережне навчальне видання}
	\texttt{Електронне навчальне видання}
	
	%{\large\today} % Date, change the \today to a set date if you want to be precise
	
	%------------------------------------------------
	%	Logo
	%------------------------------------------------
	
	%------------------------------------------------
	%	Publisher
	%------------------------------------------------
	\vspace*{4\baselineskip}
	{ \large 
		ОДЕСА\\ \vspace{-1.5mm}
		ОНУ\\ \vspace{-2.5mm}
		2026 }%Publication year
	%\vfill\vfill
	%\includegraphics[width=0.2\textwidth]{placeholder.jpg}\\[1cm] % Include a department/university logo - this will require the graphicx package
	
	%----------------------------------------------------------------------------------------
	
	
	\pagebreak
	
	\large
	\raggedright \Large {\textbf{УДК 004.4:004.421 }\\
		\hspace*{1 cm} XXX}
	
	\bigskip
	\justifying
	\large
	\hspace{1cm}\textbf{Автори:}\\
   	\textbf{\textit{Є. В. Платонова}} старший викладач кафедри оптимального керування і економічної кібернетики;\\ 
	\textbf{\textit{Є. М. Страхов}} канд. фіз.-мат. наук, доцент, доцент кафедри оптимального керування і економічної кібернетики.\\ 
	\medskip
	\hspace{1cm}\textbf{Рецензенти:}\\
	\textbf{\textit{В. О. Сперанський }}  канд. техн. наук, доцент, доцент кафедри комп’ютеризовані системи та програмні технології НУ "Одеська Політехніка";\\
	\textbf{\textit{З. Ю. Журавльова}} канд. фіз.-мат. наук, доцентом кафедри методів математичної фізики.
	
	%		\begin{tabbing}
		%			\textbf{Рецензенти: }  \= \hspace{ 1cm} \= \kill
		%			\textbf{Рецензенти:} \> \textit{\textbf{Полякова М.В.}}, докт. техн. наук, проф., професор кафедри\\ \> прикладної математики та інформаційних технологій\\ \> НУ "Одеська Політехніка".\\ 
		%			 \> \textit{\textbf{Батюк А.Є.}}, канд.техн.наук, доцент кафедри АСУ\\ \> Інституту комп'ютерних наук та інформаційних технологій\\ \> НУ "Львівська політехніка". \\
		%			 \> \textit{\textbf{Батюк А.Є.}}, канд.техн.наук, доцент кафедри \\ \> Одеського національного університету імені І. І. Мечникова.\\ \> НУ "Львівська політехніка". 
		%		\end{tabbing}
	%		\begin{tabbing}
		%				Відповідальний \= \hspace{3cm}\= \kill
		%				Відповідальний \> \\
		%				редактор: \> \textit{Тарасенко О.В.}, кандидат техн. наук, доцент
		%		\end{tabbing}
	
	
	
	\bigskip
	\medskip
	
	\centering {\normalsize \textit{Рекомендовано до видання Методичною радою ОНУ імені І.І. Мечникова\\
			(протокол № \underline{\hspace{1cm}} від \underline{\hspace{1cm}}.\underline{\hspace{1cm}}.2026 р.) } } \\
	%		за поданням Вченої ради факультету математики, фізики\\ та інформаційних технологій\\
	%		(протокол № \underline{\hspace{1cm}} від \underline{\hspace{1cm}}.\underline{\hspace{1cm}}.2024 р.)}

\bigskip
\medskip

\hspace{1.25cm} \raggedright \textbf{Платонова Є.В., Страхов Є.М.}\\ 
\textbf{XXX} \hspace{0.25cm}\justifying \indent Основи програмування на мові Python. Структурне програмування [Електронний ресурс]: конспект лекцій для здобув. першого (бакалавр) рівня вищої освіти спец. E7 "Математика" / Є.В. Платонова, Є.М. Страхов --Електронні текстові дані (1 файл: 364 Кб). Одеса: Одес. нац. ун-т ім. І. І. Мечникова, 2026. -- 55 с.
ISBN   XXX-XXX-XXX-XXX-X 

\bigskip
%\parbox{\textwidth}
\justifying 
{\textit {\small Конспект лекцій є введенням у базові принципи програмування мовою Python та систематизацією основних алгоритмічних конструкцій. У матеріалах конспекту розглядаються три базові конструкції управління виконанням програми, надаються приклади реалізації алгоритмів у Python. 
	Конспект лекцій розроблений для студентів першого курсу спеціальності E7 «Математика» та може бути використаний для самостійного навчання. Матеріал спрямований на формування 
	алгоритмічного мислення та набуття практичних навичок програмування мовою Python.}

\raggedleft	\textbf{УДК 004.4:004.421 }

\small
\begin{tabbing}
	ISBN   XXX-XXX-XXX-XXX-X \= \hspace{4cm} \= ©  \= Платонова Є.В., \= Страхов Є.М., 2026 \\
	\>   \hspace{3cm} ©  Одеський національний університет\\  \> \> \> імені І. І. Мечникова, 2026
\end{tabbing}
%	ISBN   XXX-XXX-XXX-XXX-X                                                     ©  Платонова Є. В., 2026
%																										©  Одеський національний університет	імені І. І. Мечникова, 2025
\end{titlepage}

\linespread{1.25}

\KOMAoptions{fontsize=14pt} 

